\subsection{Monte Carlo Foundations}
% Random sampling, estimators, variance, confidence intervals, law of large numbers in computation.

Monte Carlo methods turn integrals into averages of simulated random variables. This is especially useful when the integral is high-dimensional or when the target density is known only up to a normalizing constant.

\begin{method} \label{method:basic-monte-carlo-integration} \hlt{Basic Monte Carlo Integration}
\begin{enumerate}[label=(\roman*)]
\setlength{\itemsep}{0pt}
\item Let $I=\int h(x)f(x)\,dx$, where $f$ is a probability density.
\item For IID draws $X_1,\ldots,X_N\sim f$, define $\widehat{I}_N=N^{-1}\sum\limits_{i=1}^{N}h(X_i)$; then $\widehat{I}_N\xrightarrow{p}I$ by the law of large numbers.
\item With $Y_i=h(X_i)$, set $s^2=(N-1)^{-1}\sum\limits_{i=1}^{N}(Y_i-\widehat{I}_N)^2$ and $\widehat{\operatorname{se}}=\frac{s}{\sqrt{N}}$.
\item An approximate $1-\alpha$ Monte Carlo interval is $\widehat{I}_N\pm z_{\frac{\alpha}{2}}\widehat{\operatorname{se}}$.
\end{enumerate}
\end{method}

\begin{example} \label{ex:monte-carlo-integration-lln} \hlt{Monte Carlo Integration from the LLN}\\
Let $X_1,\ldots,X_n$ be iid $\operatorname{Uniform}(0,1)$ and $f\in L^1([0,1])$.\\
The estimator $\widehat I_n(f)\defeq n^{-1}\sum\limits_{i=1}^{n}f(X_i)$ satisfies $\widehat I_n(f)\xrightarrow{\mathbb{P}}\mathbb{E}f(X_1)=\int_0^1f(x)\,dx$.\\
For $f(x)=(2\pi)^{-\frac{1}{2}}e^{-\frac{x^2}{2}}$ on $[0,1]$, a $1000$-point estimate near $0.3417$ is consistent with $\int_0^1f(x)\,dx\approx0.3413$.
\end{example}

\begin{example} \label{ex:uniform-monte-carlo-integral} \hlt{Integrating on an Interval}\\
For $X\sim\operatorname{Uniform}(a,b)$, $I=\int_a^b h(x)\,dx=(b-a)\mathbb{E}[h(X)]$.\\
Given IID draws $X_1,\ldots,X_N\sim\operatorname{Uniform}(a,b)$, estimate $I$ by $\widehat{I}_N=(b-a)N^{-1}\sum\limits_{i=1}^{N}h(X_i)$.
\end{example}

\begin{example} \label{ex:monte-carlo-normal-cdf} \hlt{Estimating a Normal CDF}\\
Let $Z\sim N(0,1)$ and $\Phi(x)=\mathbb{P}(Z\leq x)$. For IID draws $Z_1,\ldots,Z_N\sim N(0,1)$, use $h(s)=\mathbf{1}\{s\leq x\}$ and $\widehat{\Phi}(x)=N^{-1}\sum\limits_{i=1}^{N}\mathbf{1}\{Z_i\leq x\}$. Thus $\widehat{\Phi}(x)$ is the empirical fraction of draws not exceeding $x$.
\end{example}

\begin{method} \label{meth:inverse-transform-simulation-examples} \hlt{Inverse-Transform Generation}
\begin{enumerate}[label=(\roman*)]
\setlength{\itemsep}{0pt}
\item \hlt{Exponential}: if $U\sim\operatorname{Uniform}(0,1)$ and $\lambda>0$, then $T=-\lambda^{-1}\log(1-U)\sim\operatorname{Exp}(\lambda)$. (or $T=-\lambda^{-1}\log U$)
\item \hlt{Weibull}: if $F(x)=1-\exp(-(\frac{x}{\alpha})^\beta)$ for $x\geq0$ and $\alpha,\beta>0$,\\
then $F^{-1}(u)=\alpha[-\log(1-u)]^{\frac{1}{\beta}}$ and $W=\alpha[-\log(1-U)]^{\frac{1}{\beta}}\sim F$.
\item \hlt{Integer-valued inverse CDF}: if $F$ is integer-valued and $U\sim\operatorname{Uniform}(0,1)$, set $Y=\inf\{k\in\Z:F(k)\geq U\}$.\\
Equivalently $F(k-1)<U\leq F(k)$. Then $F_Y=F$.
\end{enumerate}
\hlt{Key Insight}: The simulation problem is reduced to solving $u=F(x)$ for $x$.
\end{method}

\begin{method} \label{method:rejection-sampling} \hlt{Rejection Sampling}\\
Let $f$ be a target density on an interval $[a,b]$, possibly with infinite endpoints.\\
Choose an envelope $M$ such that $M(x)\geq f(x)$ on $[a,b]$ and $A\defeq\int_a^b M(t)\,dt<\infty$. Let $m(x)\defeq \frac{M(x)}{A}$.
\begin{enumerate}[label=(\roman*)]
\setlength{\itemsep}{0pt}
\item Generate $T$ with density $m$.
\item Generate $U\sim\operatorname{Uniform}(0,1)$ independently of $T$.
\item Accept $X=T$ if $U\leq \frac{f(T)}{M(T)}$; otherwise reject $T$ and repeat.
\end{enumerate}
\hlt{Key Insight}: Sample under an easy envelope and keep only points that fall below the target curve.
\end{method}

\begin{figure}[H]
  \centering
  % Source integration: schematic redraw of the rejection-method envelope from the
% attached joint-distribution chapter.
\begin{tikzpicture}[
  x=1.0cm,
  y=1.35cm,
  >=Stealth,
  line cap=round,
  line join=round,
  every node/.style={font=\scriptsize},
  axis/.style={->, gray!70, line width=0.45pt},
  tick/.style={gray!70, line width=0.4pt},
  envelope/.style={black!70, line width=0.75pt},
  target/.style={blue!70!black, line width=0.95pt},
  accept/.style={fill=blue!12, draw=none},
  accept line/.style={blue!70!black, line width=0.85pt},
  reject line/.style={gray!65, dashed, line width=0.55pt},
  point/.style={circle, fill=blue!70!black, inner sep=0pt, minimum size=2.8pt},
  label/.style={text=black!78}
]
  \def\a{-2}
  \def\b{2}
  \def\M{1.18}
  \def\T{0.65}
  \def\U{\fT}
  \pgfmathsetmacro{\fT}{0.82*exp(-\T*\T/1.35)+0.08}

  \draw[axis] (-2.45,0) -- (2.55,0) node[black!78,right] {$x$};
  \draw[axis] (-2.2,-0.04) -- (-2.2,1.38) node[black!78,above] {$y$};

  \path[accept, domain=-2:2, samples=80, variable=\x]
    (-2,0) -- plot ({\x},{0.82*exp(-\x*\x/1.35)+0.08}) -- (2,0) -- cycle;
  \draw[envelope] (\a,0) rectangle (\b,\M);
  \node[label, above=2pt] at (0,\M) {$M$};
  \draw[target, domain=-2:2, samples=100, variable=\x]
    plot ({\x},{0.82*exp(-\x*\x/1.35)+0.08}) node[right] {$f$};

  \draw[accept line] (\T,0) -- (\T,\fT);
  \draw[reject line] (\T,\fT) -- (\T,\M);
  \node[point] at (\T,\U) {};

  \node[blue!70!black] at (0.2,0.25) {accept};
  \node[black!55] at (1.2,0.94) {reject};

  \draw[tick] (\a,0.025) -- (\a,-0.025) node[black!78,below=2pt] {$a$};
  \draw[tick] (\b,0.025) -- (\b,-0.025) node[black!78,below=2pt] {$b$};
  \draw[tick] (\T,0.025) -- (\T,-0.025) node[black!78,below=2pt] {$T$};
\end{tikzpicture}

  \caption{Rejection sample accepts when the vertical draw lies below $f(T)$.}
  \label{fig:rejection-sampling-envelope}
\end{figure}

\begin{proposition} \label{prop:rejection-sampling-correctness} \hlt{Rejection Sampling Correctness}\\
Let $f$ be a density, let $M\geq f$ with $A\defeq\int M(x)\,dx<\infty$, and set $m\defeq \frac{M}{A}$. If $T\sim m$, $U\sim\operatorname{Uniform}(0,1)$ independently, and $\mathcal{A}\defeq\{U\leq \frac{f(T)}{M(T)}\}$, then $\mathbb{P}(\mathcal{A})=\frac{1}{A}$ and $T\mid\mathcal{A}$ has density $f$.
\begin{proof} For measurable $B$, $\mathbb{P}(T\in B,\mathcal{A})=\int_B \frac{m(x)f(x)}{M(x)}\,dx=A^{-1}\int_B f(x)\,dx$.\\
Taking $B=\mathbb{R}$ gives $\mathbb{P}(\mathcal{A})=A^{-1}$, and hence $\mathbb{P}(T\in B\mid\mathcal{A})=\int_B f(x)\,dx$.\\
\hlt{Structural Commentary}: Acceptance reweights the envelope $m$ back to $f$.\\
\hlt{Key Insight}: Conditional on acceptance, the normalized accepted density is exactly the target.
\end{proof}
\end{proposition}
